\section*{Bab \ref{ch:g7:motion-average-speed} --- Grafik Gerak dan Kelajuan Rata-Rata}

\begin{solution}{exo:g7:motion-average-speed:1}
$v = d/t$; \omterm{def:g6:measurement-in-science:measuring}{satuannya} \unit{km/h} dan \unit{m/s}. Jaraknya:
$d = v \times t$; waktunya: $t = d/v$.
\end{solution}

\begin{solution}{exo:g7:motion-average-speed:2}
Feri: $45 \div 3 = \qty{15}{km/h}$. Kelinci: $100 \div 8 =
\qty{12.5}{m/s}$.
\end{solution}

\begin{solution}{exo:g7:motion-average-speed:3}
$5 \times 3.6 = \qty{18}{km/h}$; \ $72 \div 3.6 = \qty{20}{m/s}$.
\end{solution}

\begin{solution}{exo:g7:motion-average-speed:4}
Garis lurus: \omterm{def:g6:describing-motion:uniform}{gerak beraturan}; lebih curam: lebih cepat; mendatar:
berhenti. Belokan yang mendatar sedikit demi sedikit bercerita tentang
perlambatan.
\end{solution}

\begin{solution}{exo:g7:motion-average-speed:5}
Dari menit $35$ sampai $60$ --- bentangan yang ketiga: ia yang paling
tidak curam di antara bentangan yang naik, terbaca langsung dari
kecondongannya yang lebih landai.
\end{solution}

\begin{solution}{exo:g7:motion-average-speed:6}
Beraturan --- satu garis lurus. Ia memperoleh \qty{2}{km} tiap
setengah \omterm{ex:g2:measuring-time:clock}{jam}: $v = \qty{4}{km/h}$.
\end{solution}

\begin{solution}{exo:g7:motion-average-speed:7}
Jarak total dibagi waktu total, termasuk perhentiannya. Seluruh
pendakiannya: $12 \div 4 = \qty{3}{km/h}$. Berjalannya saja: $12 \div 3 =
\qty{4}{km/h}$.
\end{solution}

\begin{solution}{exo:g7:motion-average-speed:8}
Atas jarak yang sama, bentangan yang lambat sekadar berlangsung lebih
lama --- lebih banyak \omterm{ex:g2:measuring-time:clock}{jam} perjalanannya dihabiskan hidup pada langkah
yang lambat, sehingga langkah yang lambat berbicara dengan lebih
banyak suara perjalanannya.
\end{solution}

\begin{solution}{exo:g7:motion-average-speed:9}
Tebakan titik tengahnya: \qty{22.5}{km/h}. Dengan jujur: perginya $30 \div 15 = 2$
\omterm{ex:g2:measuring-time:clock}{jam}, pulangnya $30 \div 30 = 1$ \omterm{ex:g2:measuring-time:clock}{jam}; totalnya \qty{60}{km} dalam $3$
\omterm{ex:g2:measuring-time:clock}{jam}: \qty{20}{km/h}. Separuh yang lambat memiliki dua dari ketiga
\omterm{ex:g2:measuring-time:clock}{jamnya}.
\end{solution}

\begin{solution}{exo:g7:motion-average-speed:10}
Grafiknya: lurus sampai (\qty{30}{min}, \qty{20}{km}); mendatar sampai
menit $45$; lurus dan lebih curam sampai (\qty{60}{min}, \qty{35}{km}).
Rata-ratanya: \qty{35}{km} dalam satu \omterm{ex:g2:measuring-time:clock}{jam} --- \qty{35}{km/h}.
\end{solution}

\begin{solution}{exo:g7:motion-average-speed:11}
$d = 340 \times 6 = \qty{2040}{m}$ --- kira-kira \qty{2}{km}. Aturan
lamanya: $6 \div 3 = 2$ kilometer --- sepakat, karena
\qty{340}{m/s} berarti hampir tepat sekilometer tiap tiga
sekon.
\end{solution}

\begin{solution}{exo:g7:motion-average-speed:12}
Separuh pertama: \qty{30}{km} pada \qty{30}{km/h} memakan tepat $1$
\omterm{ex:g2:measuring-time:clock}{jam}. Rata-rata keseluruhan \qty{60}{km/h} atas \qty{60}{km} mengizinkan
total tepat $1$ \omterm{ex:g2:measuring-time:clock}{jam} --- yang sudah habis dibelanjakan sampai sekon
yang terakhir. Sisa \qty{30}{km}-nya harus memakan $0$ \omterm{ex:g2:measuring-time:clock}{jam}: \omterm{def:g7:motion-average-speed:formula}{kelajuan}
apa pun, sepahlawan apa pun, tidak mencapai rata-ratanya; anggarannya
sudah habis.
\end{solution}

\begin{solution}{pb:g7:motion-average-speed:1}
\textbf{1.} Perjalanan yang sigap ($30$ menit), perhentian $15$ menit
(pengantaran), perjalanan yang lebih landai ($30$ menit), lalu
perhentian $15$ menit yang kedua --- dua kayuhan, dua panggilan.
\textbf{2.} \qty{9}{km} dalam setengah \omterm{ex:g2:measuring-time:clock}{jam}: \qty{18}{km/h}.
\textbf{3.} $15 - 9 = \qty{6}{km}$ dalam setengah \omterm{ex:g2:measuring-time:clock}{jam}:
\qty{12}{km/h}.
\textbf{4.} \qty{15}{km} dalam $1.5$ \omterm{ex:g2:measuring-time:clock}{jam}: \qty{10}{km/h}.
\textbf{5.} $t = 12 \div 24 = 0.5$ \omterm{ex:g2:measuring-time:clock}{jam}: $30$ menit.
\textbf{6.} $20$ menit adalah sepertiga \omterm{ex:g2:measuring-time:clock}{jam}: $d = 18 \times
\tfrac{1}{3} = \qty{6}{km}$.
\textbf{7.} $25$ menit adalah $\tfrac{25}{60}$ \omterm{ex:g2:measuring-time:clock}{jam}; $v = 10
\div \tfrac{25}{60} = \qty{24}{km/h}$.
\textbf{8.} Jaraknya: $12 + 6 + 10 = \qty{28}{km}$; waktunya:
$30 + 20 + 25 = 75$ menit $= 1.25$ \omterm{ex:g2:measuring-time:clock}{jam}; rata-rata berputarnya:
$28 \div 1.25 = \qty{22.4}{km/h}$.
\textbf{9.} $37 \div 3 \approx \qty{12.3}{km/h}$ --- jauh di bawah
$20$: tidak ada bonus (perhentiannya kerja yang jujur, tetapi
aturannya menghitungnya).
\textbf{10.} Max: separuh yang cepat $12 \div 30 = 0.4$ \omterm{ex:g2:measuring-time:clock}{jam}; separuh
yang lambat $12 \div 15 = 0.8$ \omterm{ex:g2:measuring-time:clock}{jam}; totalnya \qty{24}{km} dalam $1.2$
\omterm{ex:g2:measuring-time:clock}{jam}: tepat \qty{20}{km/h} --- bukan \emph{di atas} $20$: tidak ada
bonus.
\textbf{11.} Angka $22.5$-nya merata-ratakan kedua \emph{bilangannya},
seakan-akan tiap langkahnya memiliki separuh \omterm{ex:g2:measuring-time:clock}{jamnya}. Sebenarnya
separuh yang lambat memiliki $0.8$ dari $1.2$ \omterm{ex:g2:measuring-time:clock}{jamnya} --- dua pertiga
suara perjalanannya --- lalu menyeret rata-rata yang sebenarnya turun
ke $20$.
\textbf{12.} Misalnya: bayarlah bonusnya menurut \omterm{def:g7:motion-average-speed:average}{kelajuan rata-rata}
\emph{berputar} --- jarak total dibagi waktu mengayuh saja, dengan
bentangan mendatar grafiknya dikecualikan. Ia \omterm{def:g6:measurement-in-science:measuring}{mengukur} langkah yang
sungguh dijaga rodanya, sehingga mengganjar kayuhan yang cepat
alih-alih pengantaran yang dilewati.
\textbf{13.} ``Bentangan yang lurus adalah langkah yang mantap;
bentangan yang mendatar adalah perhentian dengan \omterm{ex:g2:measuring-time:clock}{jamnya} masih
berjalan; bentangan yang lebih curam adalah langkah yang lebih cepat
--- kecuraman adalah \omterm{def:g7:motion-average-speed:formula}{kelajuan}. Dan jangan pernah merata-ratakan
bilangan \omterm{def:g7:motion-average-speed:formula}{kelajuan}: rata-ratakanlah perjalanannya --- jarak total dibagi
waktu total --- atau \omterm{ex:g2:measuring-time:clock}{jam} yang lambat akan mempermainkan lembarnya.''
\end{solution}
